%Daten zur Institution

%\input{Dozentdaten}

%\renewcommand{\fachbereich}{Fachbereich}

%\renewcommand{\dozent}{Prof. Dr. . }

%Klausurdaten

\renewcommand{\klausurgebiet}{ }

\renewcommand{\klausurtyp}{ }

\renewcommand{\klausurdatum}{ . 20}

\klausurvorspann {\fachbereich} {\klausurdatum} {\dozent} {\klausurgebiet} {\klausurtyp}

%Daten für folgende Punktetabelle

\renewcommand{\aeins}{  3  }

\renewcommand{\azwei}{  3   }

\renewcommand{\adrei}{  0  }

\renewcommand{\avier}{  0  }

\renewcommand{\afuenf}{  6  }

\renewcommand{\asechs}{  6  }

\renewcommand{\asieben}{  0  }

\renewcommand{\aacht}{  6  }

\renewcommand{\aneun}{  11  }

\renewcommand{\azehn}{  5  }

\renewcommand{\aelf}{  5  }

\renewcommand{\azwoelf}{  3  }

\renewcommand{\adreizehn}{  7   }

\renewcommand{\avierzehn}{  2  }

\renewcommand{\afuenfzehn}{  0  }

\renewcommand{\asechzehn}{  6  }

\renewcommand{\asiebzehn}{  63  }

\renewcommand{\aachtzehn}{    }

\renewcommand{\aneunzehn}{    }

\renewcommand{\azwanzig}{    }

\renewcommand{\aeinundzwanzig}{    }

\renewcommand{\azweiundzwanzig}{    }

\renewcommand{\adreiundzwanzig}{    }

\renewcommand{\avierundzwanzig}{    }

\renewcommand{\afuenfundzwanzig}{    }

\renewcommand{\asechsundzwanzig}{    }

\punktetabellesechzehn

\klausurnote

\newpage

\setcounter{section}{K}

__NOEDITSECTION__

\inputaufgabeklausurloesung
{3}

{

Definiere die folgenden
\zusatzklammer {kursiv gedruckten} {} {} Begriffe.
\aufzaehlungsechs{Eine
\stichwort {geodätische Kurve} {}
auf einer differenzierbaren Hyperfläche

\mavergleichskette
{\vergleichskette
{ Y
}
{ \subseteq }{ \R^n
}
{  }{
}
{    }{
}
{   }{
}
}
{}{}{.}

}{Die \stichwort {Rotationsmenge} {}
\zusatzklammer {um die $x$-Achse} {} {}
zu einer Teilmenge

\mavergleichskette
{\vergleichskette
{ T
}
{ \subseteq }{  \R \times \R_{\geq 0}
}
{  }{
}
{    }{
}
{   }{
}
}
{}{}{.}

}{Die
\stichwort {Tangentialabbildung} {}
\maabbdisp {T (\varphi)} {TM} {TN
} {}
zu einer
\definitionsverweis {differenzierbaren Abbildung}{}{}
\maabbdisp {\varphi} { M } { N
} {}
zwischen
\definitionsverweis {differenzierbaren Mannigfaltigkeiten}{}{}
\mathkor {} {M} {und} {N} {.}

}{Ein
\stichwort {orientierungstreuer} {}
Kartenwechsel auf einer
\definitionsverweis {differenzierbaren Mannigfaltigkeit}{}{}
$M$.

}{Die \stichwort {zurückgezogene Differentialform} {} $\varphi^* \omega$ zu einer Differentialform
\mathl{\omega \in
{ \mathcal E }^{ k } (  M   )}{} bezüglich einer stetig differenzierbaren Abbildung
\maabb {\varphi} { L } { M
} {}
zwischen zwei differenzierbaren Mannigfaltigkeiten
\mathkor {} {L} {und} {M} {.}

}{Ein
\stichwort {metrischer} {}
linearer Zusammenhang auf dem Tangentialbündel einer riemannschen Mannigfaltigkeit $M$.
}

}
{

\aufzaehlungsechs{Eine
\zusatzklammer {als Abbildung nach $\R^n$} {} {}
zweimal
\definitionsverweis {stetig differenzierbare Kurve}{}{}
\maabbdisp {\gamma} {I} {Y
} {}
heißt
geodätische Kurve
auf $Y$, wenn ihre
\definitionsverweis {tangentiale Beschleunigung}{}{}
überall gleich $0$ ist.
}{Die Rotationsmenge zu $T$ ist

\mathdisp {\left\{ (x, y \cos \alpha, y \sin \alpha  ) \in \R^3  \mid   (x,y) \in T , \, \alpha \in [0, 2 \pi]       \right\}} { . }

}{Unter der Tangentialabbildung
\maabbdisp {T (\varphi)} {TM} {TN
} {}
versteht man die disjunkte Vereinigung der
\definitionsverweis {Tangentialabbildungen}{}{}
in den einzelnen Punkten, also

\mavergleichskettedisp
{\vergleichskette
{ T (\varphi)
}
{ =} { \biguplus_{P \in M} T_P(\varphi)
}
{ } {
}
{ } {
}
{ } {
}
}
{}{}{.}
}{Es seien
\maabbdisp {\alpha_1} {U_1 } { V_1
} {}
und
\maabbdisp {\alpha_2} {U_2 } { V_2
} {}
orientierte Karten von $M$. Der zugehörige
\definitionsverweis {Kartenwechsel}{}{}
\maabbdisp {\psi=\alpha_2 \circ \alpha_1^{-1}} {V_1 \cap \alpha_1 (U_1 \cap U_2)} {V_2 \cap \alpha_2 (U_1 \cap U_2 )
} {}
heißt orientierungstreu, wenn für jeden Punkt
\mathl{Q \in V_1 \cap \alpha_1 (U_1 \cap U_2)}{} das
\definitionsverweis {totale Differential}{}{}
\maabbdisp {\left(D\psi\right)_{Q}} {\R^n } {\R^n
} {}
\definitionsverweis {orientierungstreu}{}{}
ist.
}{Die zurückgezogene Differentialform
\mathl{\varphi^* \omega}{} ist für
\mathl{P\in L}{} und
\mathl{v_1 , \ldots , v_k \in T_PL}{} durch

\mavergleichskettedisp
{\vergleichskette
{  ( \varphi^* \omega )(P,v_1 , \ldots , v_k )
}
{ \defeq} { \omega(\varphi(P), T_P\varphi(v_1) , \ldots ,  T_P\varphi(v_k) )
}
{ } {
}
{ } {
}
{ } {
}
}
{}{}{}
definiert.
}{Der Zusammenhang heißt
metrisch,
wenn

\mavergleichskettedisp
{\vergleichskette
{ D_V \left\langle W , Z  \right\rangle
}
{ =} { \left\langle \nabla_V W , Z  \right\rangle  + \left\langle W ,  \nabla_VZ  \right\rangle
}
{ } {
}
{ } {
}
{ } {
}
}
{}{}{.}
für beliebige Vektorfelder $V,W ,Z$ gilt.
}

 }

__NOEDITSECTION__

\inputaufgabeklausurloesung
{3}

{

Formuliere die folgenden Sätze.
\aufzaehlungdrei{Der
\stichwort {Satz über die Lösung von gewöhnlichen Differentialgleichungen auf einer differenzierbaren Hyperfläche} {}

\mavergleichskette
{\vergleichskette
{  Y
}
{ \subseteq }{  \R^n
}
{  }{
}
{    }{
}
{   }{
}
}
{}{}{.}}{Die Formel für die Berechnung des kanonischen Volumens auf einer riemannschen Mannigfaltigkeit in einer Karte.}{Der Satz über die Partition der Eins.}

}
{

\aufzaehlungdrei{\faktsituation {Es sei

\mavergleichskette
{\vergleichskette
{ W
}
{ \subseteq }{ \R^n
}
{  }{
}
{    }{
}
{   }{
}
}
{}{}{}
\definitionsverweis {offen}{}{,}
\maabb {h} { W } { \R
} {}
eine
\definitionsverweis {stetig differenzierbare Funktion}{}{}
und

\mavergleichskette
{\vergleichskette
{ Y
}
{ = }{ h^{-1}(c)
}
{  }{
}
{    }{
}
{   }{
}
}
{}{}{}
die
\definitionsverweis {Faser}{}{}
zu

\mavergleichskette
{\vergleichskette
{ c
}
{ \in }{ \R
}
{  }{
}
{    }{
}
{   }{
}
}
{}{}{,}
wobei $h$ in jedem Punkt von $Y$
\definitionsverweis {regulär}{}{}
sei.
Es sei

\mavergleichskette
{\vergleichskette
{ Y
}
{ \subseteq }{ U
}
{ \subseteq }{ \R^n
}
{    }{
}
{   }{
}
}
{}{}{}
\definitionsverweis {offen}{}{}
und sei $F$ ein differenzierbares
\definitionsverweis {Vektorfeld}{}{}
auf $U$}
\faktvoraussetzung {mit der Eigenschaft, dass für jeden Punkt

\mavergleichskette
{\vergleichskette
{ P
}
{ \in }{ Y
}
{  }{
}
{    }{
}
{   }{
}
}
{}{}{}
der Vektor $F(P)$ zum
\definitionsverweis {Tangentialraum an die Faser}{}{}
in $P$ an $Y$ gehört.}
\faktfolgerung {Dann verläuft die Lösung $v(t)$ zum
\definitionsverweis {Anfangswertproblem}{}{}
zu $F$ mit der Anfangsbedingung

\mavergleichskette
{\vergleichskette
{ v(t_0)
}
{ \in }{ Y
}
{  }{
}
{    }{
}
{   }{
}
}
{}{}{}
ganz in $Y$.}
\faktzusatz {}
\faktzusatz {}}{Es sei $M$ eine
\definitionsverweis {orientierte}{}{}
\definitionsverweis {riemannsche Mannigfaltigkeit}{}{}
und $\omega$ die
\definitionsverweis {kanonische Volumenform}{}{.}
Es sei
\maabbdisp {\alpha} { U } { V
} {}
eine
\definitionsverweis {orientierte Karte}{}{}
mit

\mavergleichskettedisp
{\vergleichskette
{ V
}
{ \subseteq} { \R^n
}
{ } {
}
{ } {
}
{ } {
}
}
{}{}{}
offen mit Koordinaten
\mathl{x_1 , \ldots , x_n}{} mit der
\definitionsverweis {metrischen Fundamentalmatrix}{}{}

\mathl{G=(g_{ij})_{1 \leq i,j \leq n}}{} und
\mathl{g= \det  G}{.} Dann ist

\mavergleichskettedisp
{\vergleichskette
{  \alpha_* \omega
}
{ =} { \sqrt{ g} dx_1 \wedge \ldots \wedge dx_n
}
{ } {
}
{ } {
}
{ } {
}
}
{}{}{.}
Für eine
\definitionsverweis {messbare Teilmenge}{}{}

\mavergleichskette
{\vergleichskette
{ T
}
{ \subseteq }{ U
}
{  }{
}
{    }{
}
{   }{
}
}
{}{}{}
ist

\mavergleichskettedisp
{\vergleichskette
{
\int_{  T   }  \omega
}
{ =} {
\int_{ \alpha(T)  }   \sqrt{g} dx_1 \wedge \ldots \wedge dx_n
}
{ =} { \int_{ \alpha(T)  }   \sqrt{g}    \,  d \lambda^n
}
{ } {
}
{ } {
}
}
{}{}{.}}{Es sei $M$ eine
\definitionsverweis {differenzierbare Mannigfaltigkeit}{}{}
mit einer
\definitionsverweis {abzählbaren Basis}{}{}
der Topologie. Dann gibt es zu jeder offenen Überdeckung eine der Überdeckung untergeordnete
\definitionsverweis {stetig differenzierbare}{}{}
\definitionsverweis {Partition der Eins}{}{.}}

 }

__NOEDITSECTION__

\inputaufgabeklausurloesung
{0}

{

}
{  }

__NOEDITSECTION__

\inputaufgabeklausurloesung
{0}

{

}
{  }

__NOEDITSECTION__

\inputaufgabeklausurloesung
{6}

{

Bestimme für die durch

\mavergleichskettedisp
{\vergleichskette
{ x^2+y^2+2z^2
}
{ =} { 4
}
{ } {
}
{ } {
}
{ } {
}
}
{}{}{}
gegebene Fläche $Y$ und den Punkt

\mavergleichskettedisp
{\vergleichskette
{ P
}
{ =} { \left(  1  , \,   1  , \,   1                 \right)
}
{ \in} { Y
}
{ } {
}
{ } {
}
}
{}{}{}
eine Diagonalmatrix für die
\definitionsverweis {Weingartenabbildung}{}{}
$L_P$.

}
{

Das normierte Gradientenfeld ist

\mavergleichskettealign
{\vergleichskettealign
{ N(x,y,z)
}
{ =} { { \frac{   1  }{  \sqrt{ 4x^2+4y^2+16 z^2}  } } \begin{pmatrix}  2x \\ 2y\\  4z   \end{pmatrix}
}
{ =} { { \frac{   1  }{  \sqrt{ x^2+y^2+4 z^2}  } } \begin{pmatrix}  x  \\ y \\  2z   \end{pmatrix}
}
{ } {
}
{ } {
}
}
{}
{}{.}
Wir arbeiten mit dem Einheitsnormalenfeld

\mavergleichskettedisp
{\vergleichskette
{ M(x,y,z)
}
{ =} { { \frac{   1  }{  \sqrt{ 4+2 z^2}  } } \begin{pmatrix}  x  \\ y \\  2z   \end{pmatrix}
}
{ } {
}
{ } {
}
{ } {
}
}
{}{}{,}
was auf $Y$ mit $N$ übereinstimmt. Das totale Differential von $M$ ist

\mathdisp {\begin{pmatrix}  { \frac{   1  }{  \sqrt{4+2z^2}  } }   &   0 &  { \frac{   -2 xz }{  \left(  4+2z^2  \right)^{3/2}  } }  \\  0  &  { \frac{   1  }{  \sqrt{4+2z^2}  } }   &  { \frac{   -2 yz }{  \left(  4+2z^2  \right)^{3/2}  } }   \\ 0   &  0  &  { \frac{   8 }{  \left(  4+2z^2  \right)^{3/2}  } }  \end{pmatrix}} { . }

Im angegebenen Punkt $P$ ist der Gradient $\begin{pmatrix}  2  \\ 2\\  4   \end{pmatrix}$ und die beiden Vektoren
\mathkor {} {\begin{pmatrix}  1  \\ -1\\  0   \end{pmatrix}} {und} {\begin{pmatrix}  2  \\ 0 \\  -1   \end{pmatrix}} {}
ist eine Basis des Tangetialraumes $T_PY$. Das totale Differential zu $M$ ist in diesem Punkt gleich

\mathdisp {\begin{pmatrix}  { \frac{   1  }{  \sqrt{6}  } }   &   0 &  { \frac{   -1 }{  3 \sqrt{6}  } }  \\  0  &  { \frac{   1  }{  \sqrt{6}  } }   &  { \frac{   -1  }{  3 \sqrt{6}  } }   \\ 0   &  0  &  { \frac{   4  }{  3 \sqrt{6}  } }   \end{pmatrix}} { . }

Angewendet auf den ersten Basisvektor $\begin{pmatrix}  1  \\ -1\\  0   \end{pmatrix}$  ergibt sich $\begin{pmatrix}  { \frac{   1  }{  \sqrt{6}  } }  \\ - { \frac{   1  }{  \sqrt{6}  } } \\  0   \end{pmatrix}$, dies ist also ein Eigenvektor zum Eigenwert ${ \frac{   1  }{  \sqrt{6}  } }$. Angewendet auf den zweiten Basisvektor $\begin{pmatrix}  2  \\ 0 \\  -1   \end{pmatrix}$  ergibt sich

\mavergleichskettedisp
{\vergleichskette
{ \begin{pmatrix}  { \frac{   7 }{  3 \sqrt{6}  } }  \\ { \frac{   1 }{  3 \sqrt{6 } }}  \\  - { \frac{   4  }{  3 \sqrt{6}  } }   \end{pmatrix}
}
{ =} { - { \frac{   1 }{  3 \sqrt{6 } }} \begin{pmatrix}  1  \\ -1\\  0   \end{pmatrix} +   { \frac{   4  }{  3 \sqrt{6}  } } \begin{pmatrix}  2  \\ 0 \\  -1   \end{pmatrix}
}
{ } {
}
{ } {
}
{ } {
}
}
{}{}{.}
Daher ist der andere Eigenwert gleich ${ \frac{   4  }{  3 \sqrt{6}  } }$ und eine beschreibende Diagonalmatrix ist

\mathdisp {\begin{pmatrix}  { \frac{   1  }{  \sqrt{6}  } }   &   0  \\  0  &  { \frac{   4  }{  3 \sqrt{6}  } }  \end{pmatrix}} { . }

 }

__NOEDITSECTION__

\inputaufgabeklausurloesung
{6}

{

Beweise den Isometriesatz für den Paralleltransport auf einer Hyperfläche

\mavergleichskette
{\vergleichskette
{ Y
}
{ \subseteq }{ \R^n
}
{  }{
}
{    }{
}
{   }{
}
}
{}{}{.}

}
{

Zum Nachweis der Linearität seien

\mavergleichskette
{\vergleichskette
{ v,w
}
{ \in }{ T_PY
}
{  }{
}
{    }{
}
{   }{
}
}
{}{}{}
and

\mavergleichskette
{\vergleichskette
{ r,s
}
{ \in }{ \R
}
{  }{
}
{    }{
}
{   }{
}
}
{}{}{}
gegeben. Es seien
\mathkor {} {F} {bzw.} {G} {}
die gemäß
[[Hyperfläche/Kurve/Anfangstangentenvektor/Paralleles Vektorfeld/Existenz/Fakt|Satz 6.9 (Differentialgeometrie (Osnabrück 2023))]]
eindeutig bestimmten parallelen Vektorfelder längs $\gamma$ mit

\mavergleichskette
{\vergleichskette
{ F(a)
}
{ = }{ v
}
{  }{
}
{    }{
}
{   }{
}
}
{}{}{}
und

\mavergleichskette
{\vergleichskette
{ G(a)
}
{ = }{ w
}
{  }{
}
{    }{
}
{   }{
}
}
{}{}{.}
Nach
[[Hyperfläche/Kurve/Paralleles Vektorfeld/Eigenschaften/Fakt|Lemma 6.7 (Differentialgeometrie (Osnabrück 2023))]]
ist
\mathl{rF+sG}{} ein paralleles Vektorfeld mit

\mavergleichskette
{\vergleichskette
{ (rF+sG)(a)
}
{ = }{ v+ w
}
{  }{
}
{    }{
}
{   }{
}
}
{}{}{.}
Wegen der Eindeutigkeit aus
[[Hyperfläche/Kurve/Anfangstangentenvektor/Paralleles Vektorfeld/Existenz/Fakt|Satz 6.9 (Differentialgeometrie (Osnabrück 2023))]]
ist somit
\mathl{rF+sG}{} das parallele Vektorfeld zum Tangentialvektor
\mathl{rv+sw}{.} Daher ist

\mavergleichskettedisp
{\vergleichskette
{ \Psi_{\gamma} (rv+sw)
}
{ =} { (rF+sG)(b)
}
{ =} { r F(b) +s G(b)
}
{ =} { r \Psi_{\gamma} (v) + s \Psi_{\gamma} (w)
}
{ } {
}
}
{}{}{.}

Zum Nachweis der Verträglichkeit mit dem Skalarprodukt seien wieder

\mavergleichskette
{\vergleichskette
{ v,w
}
{ \in }{ T_PY
}
{  }{
}
{    }{
}
{   }{
}
}
{}{}{}
gegeben und es seien $F,G$ die zugehörigen parallelen Vektorfelder. Es ist

\mavergleichskettedisp
{\vergleichskette
{ \left\langle  F(t) , G(t)  \right\rangle'
}
{ =} { \left\langle  F'(t) , G(t)  \right\rangle  +  \left\langle  F(t) , G'(t)  \right\rangle
}
{ =} { 0
}
{ } {
}
{ } {
}
}
{}{}{,}
da $F,G$ tangential sind und $F',G'$ orthogonal zum Tangentialraum sind. Daher ist
\mathl{\left\langle  F(t) , G(t)  \right\rangle}{} konstant längs des Weges. Daher ist

\mavergleichskettedisp
{\vergleichskette
{ \left\langle \Psi_\gamma(v)  , \Psi_\gamma(w)  \right\rangle
}
{ =} { \left\langle  F(b)  ,  G(b)   \right\rangle
}
{ =} { \left\langle  F(a)  ,  F(a)   \right\rangle
}
{ =} { \left\langle  v  , w  \right\rangle
}
{ } {
}
}
{}{}{.}
Die Bijektivität ist damit auch klar.

 }

__NOEDITSECTION__

\inputaufgabeklausurloesung
{0}

{

}
{  }

__NOEDITSECTION__

\inputaufgabeklausurloesung
{6}

{

Es sei $M$ eine kompakte topologische $d$-dimensionale Mannigfaltigkeit,

\mavergleichskette
{\vergleichskette
{ d
}
{ \geq }{ 1
}
{  }{
}
{    }{
}
{   }{
}
}
{}{}{.}
Zeige, dass es eine beschränkte offene Teilmenge

\mavergleichskette
{\vergleichskette
{ U
}
{ \subseteq }{ \R^d
}
{  }{
}
{    }{
}
{   }{
}
}
{}{}{}
und eine stetige surjektive Abbildung
\maabbdisp {\varphi} { U } { M
} {}
gibt.

}
{

Zu jedem Punkt
\mathl{P \in M}{} wählen wir eine offene Kartenumgebung
\mathl{P \in U_P}{} mit einer Karte
\maabbdisp {\alpha_P} {U_P } {V_P
} {}
mit
\mathl{V_P \subseteq \R^d}{.} Dabei können wir annehmen, indem wir zu einer Ballumgebung von
\mathl{\alpha_P(P) \in V_P}{} und dessen Urbild übergehen, dass die $V_p$ offene Bälle sind, deren Radius maximal
\mathl{{ \frac{   1 }{  3 } }}{} ist. Die

\mathbed {U_P} {}
 {P\in M} {}
 {} {} {} {,}
überdecken die Mannigfaltigkeit. Wegen der Kompaktheit von $M$ gibt es endlich viele Punkte
\mathl{P_1 , \ldots , P_n}{} derart, dass auch

\mathbed {U_i=U_{P_i}} {}
 {i=1 , \ldots , n} {}
 {} {} {} {,}
die Mannigfaltigkeit überdecken. Wir platzieren die offenen Bälle
\mathl{V_i=V_{P_i}}{} in
\zusatzklammer {\anfuehrung{einem neuen}{}} {} {}
$\R^d$, und zwar mit den Mittelpunkten

\mathdisp {(1,0 , \ldots , 0), (2,0 , \ldots , 0) , \ldots , (n,0 , \ldots , 0)} { . }

Wegen der Radiusbedingung sind diese Bälle zueinander disjunkt. Wir betrachten die beschränkte offene Menge
\mathl{U= \bigcup_{i=1}^n V_i}{.} Die Kartenabbildungen $\alpha_i$ liefern stetige Abbildungen

\mathdisp {\varphi_i :V_i \stackrel{ \left(  \alpha_i  \right)^{-1} }{\longrightarrow} U_i \longrightarrow M} { . }

Wegen der Disjunktheit ergibt sich daraus durch
\mathl{\varphi(z)= \varphi_i(z)}{,} falls
\mathl{z \in V_i}{} ist, eine stetige Abbildung
\maabbdisp {\varphi} { U } {M
} {.}
Wegen der Überdeckungseigenschaft ist diese surjektiv.

 }

__NOEDITSECTION__

\inputaufgabeklausurloesung
{11 (1+3+1+2+4)}

{

Wir betrachten die Menge

\mavergleichskettedisp
{\vergleichskette
{N
}
{ =} { \left\{ A= \begin{pmatrix}  a   &   b   \\  c  & d \end{pmatrix}   \mid   A \text{ ist nilpotent}         \right\}
}
{ \subseteq} {\R^4
}
{ } {
}
{ } {
}
}
{}{}{}
der reellen nilpotenten
$2\times 2$-\definitionsverweis {Matrizen}{}{}
sowie die Menge

\mavergleichskettedisp
{\vergleichskette
{M
}
{ =} { N \setminus \{ \begin{pmatrix}  0   &   0  \\  0  &  0 \end{pmatrix}  \}
}
{ } {
}
{ } {
}
{ } {
}
}
{}{}{.}

a) Ist $N$ zusammenhängend?

b) Zeige, dass $M$ eine
\definitionsverweis {abgeschlossene Untermannigfaltigkeit}{}{}
einer offenen Teilmenge

\mavergleichskette
{\vergleichskette
{ G
}
{ \subseteq }{ \R^4
}
{  }{
}
{    }{
}
{   }{
}
}
{}{}{}
ist.

c) Bestimme die
\definitionsverweis {Dimension}{}{}
von $M$.

d) Ist $M$ zusammenhängend?

e) Überdecke
\mathl{M}{} mit expliziten topologischen Karten.

}
{

a) Jede nilpotente Matrix
\mathl{\begin{pmatrix}  a   &   b   \\  c  & d \end{pmatrix}}{} lässt sich durch den linearen Weg
\maabbeledisp {} {[0,1]} {N
} { t } {t \begin{pmatrix}  a   &   b   \\  c  & d \end{pmatrix}
} {,}
innerhalb der nilpotenten Matrizen mit der Nullmatrix verbinden. Daher ist $N$ wegzusammenhängend und damit auch zusammenhängend.

b) Eine
$2 \times 2$-\definitionsverweis {Matrix}{}{}
ist genau dann nilpotent, wenn sowohl die Spur als auch die Determinante $0$ sind. Die Menge $N$ der nilpotenten Matrizen kann also als

\mathdisp {\left\{  (a,b,c,d)  \mid   a+d = 0 \text{ und } ad-bc = 0         \right\}} {  }

aufgefasst werden. Wir betrachten die Abbildung
\maabbeledisp {} {\R^4} {\R^2
} {(a,b,c,d)} { (a+d, ad-bc)
} {.}
Deren Jacobi-Matrix ist

\mathdisp {\begin{pmatrix}  1 &  0 &  0 &  1 \\  d  &  -c &  -b &  a  \end{pmatrix}} { . }

Diese Abbildung ist im Nullpunkt nicht
\definitionsverweis {regulär}{}{,}
aber in jedem anderen Punkt der Faser $N$. Wenn nämlich
\mathl{P \in M}{} ist, so folgt wegen

\mavergleichskettedisp
{\vergleichskette
{a^2
}
{ =} {bc
}
{ } {
}
{ } {
}
{ } {
}
}
{}{}{}
aus

\mavergleichskettedisp
{\vergleichskette
{b
}
{ =} { c
}
{ =} { 0
}
{ } {
}
{ } {
}
}
{}{}{}
sofort

\mavergleichskettedisp
{\vergleichskette
{a
}
{ =} {b
}
{ =} { c
}
{ =} { d
}
{ =} { 0
}
}
{}{}{.}
Die Jacobi-Matrix hat also in den Punkten aus $M$ den maximalen Rang. Mit

\mavergleichskettedisp
{\vergleichskette
{G
}
{ =} {\R^4 \setminus \{0\}
}
{ } {
}
{ } {
}
{ } {
}
}
{}{}{}
kann man $M$ als die Faser der eingeschränkten Abbildung auffassen, die überall auf der Faser regulär ist. Daher ist $M$ nach
[[Satz über implizite Abbildungen/Faser ist Mannigfaltigkeit/Fakt|Satz 8.2 (Differentialgeometrie (Osnabrück 2023))]]
eine abgeschlossene Untermannigfaltigkeit von $G$.

c) Nach
[[Satz über implizite Abbildungen/Faser ist Mannigfaltigkeit/Fakt|Satz 8.2 (Differentialgeometrie (Osnabrück 2023))]]
ist die Dimension von $M$ gleich

\mavergleichskette
{\vergleichskette
{ 4-2
}
{ = }{ 2
}
{  }{
}
{    }{
}
{   }{
}
}
{}{}{.}

d) Wir schreiben

\mavergleichskettedisp
{\vergleichskette
{M_+
}
{ =} { \left\{ (a,b,c,d) \in M  \mid   b>c        \right\}
}
{ } {
}
{ } {
}
{ } {
}
}
{}{}{}
und

\mavergleichskettedisp
{\vergleichskette
{M_-
}
{ =} { \left\{ (a,b,c,d) \in M  \mid   b<c        \right\}
}
{ } {
}
{ } {
}
{ } {
}
}
{}{}{,}
beides sind
\zusatzklammer {als Durchschnitt von $M$ mit der durch
\mathl{b>c}{} gegebenen offenen Menge des $\R^4$} {} {}
offene Mengen in $M$. Die Matrizen
\mathkor {} {\begin{pmatrix}  0  &   1  \\  0 &  0 \end{pmatrix}} {und} {\begin{pmatrix}  0  &   0  \\  1 &  0 \end{pmatrix}} {}
zeigen, dass sie nicht leer sind. Ferner überdecken sie ganz $M$. Bei

\mavergleichskette
{\vergleichskette
{b
}
{ = }{ c
}
{  }{
}
{    }{
}
{   }{
}
}
{}{}{}
folgt nämlich wegen

\mavergleichskettedisp
{\vergleichskette
{ 0
}
{ =} { ad-bc
}
{ =} { -a^2 -b^2
}
{ } {
}
{ } {
}
}
{}{}{}
direkt

\mavergleichskette
{\vergleichskette
{a
}
{ = }{b
}
{ = }{ c
}
{ =   }{ d
}
{ =  }{ 0
}
}
{}{}{,}
und der Punkt gehört nicht zu $M$. Es liegt also eine Überdeckung mit zwei nichtleeren disjunkten offenen Mengen vor, daher ist $M$ nicht zusammenhängend.

e) Wir arbeiten mit der Abbildung
\maabbeledisp {\psi} {\R^2 \setminus \{(0,0)\}} { M_+
} {(a,u)} { (a, u + \sqrt{a^2+u^2}, u - \sqrt{a^2 +u^2},-a)
} {.}
Wegen

\mavergleichskettealign
{\vergleichskettealign
{ad-bc
}
{ =} { -a^2 - \left(  u + \sqrt{a^2+u^2}  \right) \left( u - \sqrt{a^2+u^2}  \right)
}
{ =} { -a^2 - \left(  u^2 - \left(  a^2+u^2 \right)   \right)
}
{ =} { 0
}
{ } {
}
}
{}
{}{}
ist die Determinantenbedingung erfüllt, und wegen

\mavergleichskettedisp
{\vergleichskette
{b-c
}
{ =} { u + \sqrt{a^2+u^2} - \left(  u- \sqrt{a^2 +u^2})  \right)
}
{ =} { 2 \sqrt{a^2+u^2}
}
{ >} { 0
}
{ } {
}
}
{}{}{}
gehört das Bild zu $M_+$. Die Abbildung
\maabbeledisp {\varphi} {M_+} { \R^2 \setminus \{(0,0)\}
} {(a,b,c,d)} { \left(  a  , \,   { \frac{  b+c }{  2 } }                   \right)
} {,}
ist invers zu der gegebenen Abbildung. Dabei ist

\mavergleichskettedisp
{\vergleichskette
{ \varphi \circ \psi
}
{ =} {
\operatorname{Id}_{ \R^2 \setminus \{(0,0)\}  }
}
{ } {
}
{ } {
}
{ } {
}
}
{}{}{}
klar. Die andere Identität ergibt sich aus

\mavergleichskettealign
{\vergleichskettealign
{ { \frac{  b+c }{  2 } } + \sqrt{a^2 + \left( { \frac{  b+c }{  2 } }  \right)^2 }
}
{ =} { { \frac{  b+c }{  2 } } + { \frac{   1 }{  2 } }  \sqrt{4 a^2 + b^2 +c^2 +2bc }
}
{ =} { { \frac{  b+c }{  2 } } + { \frac{   1 }{  2 } }  \sqrt{ b^2 +c^2 -2bc }
}
{ =} { { \frac{  b+c }{  2 } } + { \frac{  b-c }{  2 } }
}
{ =} {b
}
}
{}
{}{}
und

\mavergleichskettedisp
{\vergleichskette
{ { \frac{  b+c }{  2 } } - \sqrt{a^2 + \left(   { \frac{  b+c }{  2 } }   \right)^2 }
}
{ =} { { \frac{  b+c }{  2 } } - { \frac{  b-c }{  2 } }
}
{ =} { c
}
{ } {}
{ } {}
}
{}{}{.}
Beide Abbildungen sind stetig, daher liegt eine
\definitionsverweis {Homöomorphie}{}{}
vor.
Für
\mathl{M_-}{} vertauscht man die Rollen von
\mathkor {} {b} {und} {c} {.}

 }

__NOEDITSECTION__

\inputaufgabeklausurloesung
{5}

{

Sagen Sie etwas Schlaues zum Thema Möbiusband.

}
{  }

__NOEDITSECTION__

\inputaufgabeklausurloesung
{5}

{

Wir betrachten die Abbildung
\maabbeledisp {\varphi} { \R^3 } { \R^2
} { (x,y,z) } { \left(  xyz^2,xy-z^3  \right)
} {.}
Berechne die Matrix der Abbildung
\maabbdisp {\bigwedge^2 T_P (\varphi)} { \bigwedge^2 T_P \R^3 } { \bigwedge^2 T_{\varphi(P)}  \R^2
} {}
im Punkt

\mavergleichskette
{\vergleichskette
{ P
}
{ = }{ (1,3,5)
}
{  }{
}
{    }{
}
{   }{
}
}
{}{}{}
bezüglich einer geeigneten Basis.

}
{

Die Jacobimatrix von $\varphi$ ist allgemein

\mathdisp {\begin{pmatrix}  yz^2 &  xz^2 &  2xyz \\  y  &  x  &  -3z^2 \end{pmatrix}} { . }

Für den Punkt

\mavergleichskette
{\vergleichskette
{ P
}
{ = }{ (1,3,5)
}
{  }{
}
{    }{
}
{   }{
}
}
{}{}{}
liegt daher die Jacobimatrix

\mathdisp {\begin{pmatrix}  75 &  25 &  30 \\  3 &  1 &  -75  \end{pmatrix}} {  }

vor. Diese Matrix beschreibt die lineare Abbildung
\maabbdisp {L} {T_P\R^3 \cong \R^3} { T_{\varphi(P)} \R^2 \cong \R^2
} {}
bezüglich der Standardbasen. Wir bestimmen die Matrixdarstellung für die Abbildung
\maabbdisp {\bigwedge^2 L} {\bigwedge^2 \R^3 } { \bigwedge^2 \R^2
} {}
bezüglich der Basen
$e_1 \wedge e_2,\, e_1 \wedge e_3$ , $e_2 \wedge e_3$
\zusatzklammer {links} {} {}
und
\mathl{e_1 \wedge e_2}{}
\zusatzklammer {rechts} {} {.}
Dazu müssen wir die Bilder dieser Dachprodukte ausrechnen. Es ist

\mavergleichskettealign
{\vergleichskettealign
{ (\bigwedge^2 L ) (e_1 \wedge e_2)
}
{ =} { L(e_1) \wedge L(e_2)
}
{ =} { (75e_1 + 3e_2) \wedge (  25 e_1 + 1 e_2 )
}
{ =} { 75 e_1 \wedge e_2 + 75 e_2 \wedge  e_1
}
{ =} { 75 e_1 \wedge e_2 - 75 e_1 \wedge  e_2
}
}
{
\vergleichskettefortsetzungalign
{ =} { 0
}
{ } {}
{ } {}
{ } {}
}
{}{,}

\mavergleichskettealign
{\vergleichskettealign
{ (\bigwedge^2 L ) (e_1 \wedge e_3)
}
{ =} { L(e_1) \wedge L(e_3)
}
{ =} { (75e_1 + 3e_2) \wedge (  30 e_1 - 75 e_2 )
}
{ =} { - 75^2 e_1 \wedge e_2 + 90 e_2 \wedge  e_1
}
{ =} { (-5625 - 90) e_1 \wedge e_2
}
}
{
\vergleichskettefortsetzungalign
{ =} { -5715 e_1 \wedge e_2
}
{ } {}
{ } {}
{ } {}
}
{}{,}
und

\mavergleichskettealign
{\vergleichskettealign
{ (\bigwedge^2 L ) (e_2 \wedge e_3)
}
{ =} { L(e_2) \wedge L(e_3)
}
{ =} { (  25 e_1 + 1 e_2 )   \wedge ( 30 e_1 - 75 e_2)
}
{ =} { - 1875 e_1 \wedge e_2 + 30 e_2 \wedge  e_1
}
{ =} { -  1905 e_1 \wedge  e_2
}
}
{}
{}{.}
Die beschreibende Matrix ist also

\mathdisp {\left(  0,-5715,-1905                     \right)} { . }

 }

__NOEDITSECTION__

\inputaufgabeklausurloesung
{3}

{

Es seien

\mavergleichskette
{\vergleichskette
{ W
}
{ \subseteq }{ \R^m
}
{  }{
}
{    }{
}
{   }{
}
}
{}{}{}
und

\mavergleichskette
{\vergleichskette
{ U
}
{ \subseteq }{ \R^n
}
{  }{
}
{    }{
}
{   }{
}
}
{}{}{}
\definitionsverweis {offene Teilmengen}{}{}
und sei
\maabbdisp {\psi} { W } { U
} {}
eine
\definitionsverweis {stetig differenzierbare}{}{}
\definitionsverweis {Abbildung}{}{.}
Es sei
\maabbdisp {f} { U } {\R
} {}
eine stetig differenzierbare Funktion. Folgere aus der
[[Totale Differenzierbarkeit/K/Kettenregel/Fakt|Kettenregel]],
dass

\mavergleichskettedisp
{\vergleichskette
{ d (\psi^*f)
}
{ =} { \psi^*(df)
}
{ } {
}
{ } {
}
{ } {
}
}
{}{}{}
gilt, wobei $\psi^*$ das
\definitionsverweis {Zurückziehen von Differentialformen}{}{}
bezeichnet.

}
{

Es sei

\mavergleichskette
{\vergleichskette
{ P
}
{ \in }{  W
}
{  }{
}
{    }{
}
{   }{
}
}
{}{}{}
und

\mavergleichskette
{\vergleichskette
{ v
}
{ \in }{  \R^m
}
{  }{
}
{    }{
}
{   }{
}
}
{}{}{.}
Es ist einerseits

\mavergleichskettedisp
{\vergleichskette
{ d( \psi^* f) (P,v)
}
{ =} { D_P( f \circ \psi) (v)
}
{ =} { \left(  \left(  D_{\psi(P)} f  \right)  \circ \left(  D_P  \psi  \right)  \right) (v)
}
{ =} { D_{\psi(P)} (f) \left(  D_P \psi(v)  \right)
}
{ } {
}
}
{}{}{.}
Andererseits ist auch

\mavergleichskettedisp
{\vergleichskette
{ \left(  \psi^*(df)  \right) (P,v)
}
{ =} { (df) \left(  \psi(P), (D_P \psi )(v)  \right)
}
{ =} { D_{\psi(P)} (f) \left(  D_P \psi (v)  \right)
}
{ } {
}
{ } {
}
}
{}{}{.}

 }

__NOEDITSECTION__

\inputaufgabeklausurloesung
{7}

{

Es sei $M$ eine
\definitionsverweis {differenzierbare Mannigfaltigkeit}{}{}
mit
\definitionsverweis {abzählbarer Basis der Topologie}{}{}
und sei $\omega$ eine
\definitionsverweis {positive Volumenform}{}{}
auf $M$. Es sei
\maabbdisp {\varphi} { L } { M
} {}
ein
\definitionsverweis {Diffeomorphismus}{}{}
mit der Mannigfaltigkeit $L$ und

\mavergleichskette
{\vergleichskette
{ S
}
{ \subseteq }{ L
}
{  }{
}
{    }{
}
{   }{
}
}
{}{}{}
eine
\definitionsverweis {messbare Teilmenge}{}{.}
Zeige

\mavergleichskettedisp
{\vergleichskette
{ \int_S \varphi^* \omega
}
{ =} { \int_{ \varphi(S)} \omega
}
{ } {
}
{ } {
}
{ } {
}
}
{}{}{.}

}
{

Nach
[[Mannigfaltigkeit/Abzählbar/Positive Volumenform/Zugehöriges Maß/Vorbereitende Eigenschaften/Fakt|Lemma 15.2 (Differentialgeometrie (Osnabrück 2023))]]
können wir davon ausgehen, dass
\mathkor {} {S} {und} {\varphi(S)} {}
jeweils ganz in einer Karte von
\mathkor {} {L} {bzw. von} {M} {}
liegen, sagen wir

\mavergleichskette
{\vergleichskette
{ S
}
{ \subseteq }{ U
}
{  }{
}
{    }{
}
{   }{
}
}
{}{}{}
mit der Karte
\maabbdisp {\alpha} { U } { U' \subseteq \R^n
} {}
und

\mavergleichskette
{\vergleichskette
{ \varphi(S)
}
{ \subseteq }{ V
}
{  }{
}
{    }{
}
{   }{
}
}
{}{}{}
mit der Karte
\maabbdisp {\beta} { V } { V' \subseteq \R^n
} {.}
Wir können weiter durch verkleinern der Karten davon ausgehen, dass
\mathkor {} {U} {und} {V} {}
über $\varphi$ diffeomorph zueinander sind. Es liegt dann ein kommutatives Diagramm von Diffeomorphismen

\mathdisp {\begin{matrix}  U  & \stackrel{ \varphi }{\longrightarrow} &  V  &  \\ \!\!\!\!\! \alpha \downarrow & & \downarrow \beta \!\!\!\!\!  & \\  U'  & \stackrel{ \beta \circ \varphi \circ \alpha^{-1} }{\longrightarrow} &  V' & \! \! \! \!\!\!\!\!   \\ \end{matrix}} {  }

vor, das ein kommutatives Diagramm

\mathdisp {\begin{matrix}  S  & \stackrel{ \varphi }{\longrightarrow} &  \varphi (S)  &  \\ \!\!\!\!\! \alpha \downarrow & & \downarrow \beta \!\!\!\!\!  & \\  \alpha (S) & \stackrel{ \beta \circ \varphi \circ \alpha^{-1} }{\longrightarrow} &  \beta ( \varphi (S) ) & \! \! \! \!\!\!\!\!   \\ \end{matrix}} {  }

von messbaren Mengen induziert. Für die Volumenform $\omega$ auf $V$ gilt dabei

\mavergleichskettedisp
{\vergleichskette
{ \alpha^{ {-1}*} ( \varphi^* \omega)
}
{ =} { \left(  \beta \circ \varphi \circ \alpha^{-1}  \right)^*  ( \beta^{ {-1}*} \omega)
}
{ } {
}
{ } {
}
{ } {
}
}
{}{}{.}
nach
[[Mannigfaltigkeit/Differenzierbare Abbildung/Zurückziehen von Differentialformen/Elementare Eigenschaften/Fakt|Lemma 14.7 (Differentialgeometrie (Osnabrück 2023))  (4)]].
Die Volumenform $\beta^{ {-1}*} \omega$ besitzt auf $V'$ die Beschreibung
\mathl{g dy_1  \wedge \ldots \wedge dy_n}{} mit einer messbaren Funktion
\maabb {g} {V'} {\R
} {}
und $\int_{\varphi(S)} \omega$ ist nach Definition gleich $\int_{\beta( \varphi(S))} g d \lambda^n$. Ebenso besitzt
\mathl{\alpha^{ {-1}*} ( \varphi^* \omega)}{} auf $U'$ eine Beschreibung der Form
\mathl{f dx_1 \wedge \ldots \wedge dx_n}{.} Diese Volumenform auf $U'$ stimmt mit dem Rückzug von
\mathl{g dy_1  \wedge \ldots \wedge dy_n}{} überein und dieser ist nach
[[Offene Mengen/Differenzierbare Abbildung/Zurückziehen von Volumenform auf gleichdimensionalem Raum/In Koordinaten/Fakt|Korollar 14.9 (Differentialgeometrie (Osnabrück 2023))]]
gleich

\mathdisp {(g \circ \psi) \cdot \det  \left(  \left(  { \frac{   \partial   \psi_i   }{  \partial   x_j   } }  \right)_{1 \leq  i, j \leq n}  \right) dx_1 \wedge \ldots \wedge dx_n} {  }

\zusatzklammer {mit

\mavergleichskettek
{\vergleichskettek
{ \psi
}
{ = }{ \beta \circ \varphi \circ \alpha^{-1}
}
{  }{
}
{    }{
}
{   }{
}
}
{}{}{}} {} {.}
Somit folgt die Aussage aus
[[Diffeomorphismus/Transformationsformel für Integrale/Fakt|[[:Kurs:Maß- und Integrationstheorie (Osnabrück 2022-2023)/3/Klausur mit Lösungen/latex]] (Maß- und Integrationstheorie (Osnabrück 2022-2023))[[Kategorie:Kurs:Differentialgeometrie/Referenznummer|Maß- und Integrationstheorie (Osnabrück 2022-2023)]] (Differentialgeometrie (Osnabrück 2023))]].

 }

__NOEDITSECTION__

\inputaufgabeklausurloesung
{2}

{

Wir betrachten auf dem
\definitionsverweis {trivialen Vektorbündel}{}{}
\maabbdisp {p} {\R^2 \times \R} { \R^2
} {}
vom
\definitionsverweis {Rang}{}{}
$1$ über $\R^2$ den
\definitionsverweis {linearen Zusammenhang}{}{,}
der durch die
\definitionsverweis {Christoffelsymbole}{}{}

\mavergleichskette
{\vergleichskette
{ \Gamma_1 (x,y)
}
{ = }{ x^5-x^2y^2+3y^4
}
{  }{
}
{    }{
}
{   }{
}
}
{}{}{}
und

\mavergleichskette
{\vergleichskette
{ \Gamma_2 (x,y)
}
{ = }{ 7x^3 +xy^2-4y^5
}
{  }{
}
{    }{
}
{   }{
}
}
{}{}{}
gegeben sei. Berechne den
\definitionsverweis {Krümmungsoperator}{}{}
$R \left(  \partial_1, \partial_2  \right)$.

}
{

Nach
[[Reelles Vektorbündel/Rang 1/Offene Menge/Krümmungsoperator/Fakt|Fakt]] [[Kurs:Differentialgeometrie/Reelles Vektorbündel/Rang 1/Offene Menge/Krümmungsoperator/Fakt/Faktreferenznummer|*****]]
ist

\mavergleichskettealign
{\vergleichskettealign
{ R \left(  \partial_1, \partial_2  \right)
}
{ =} { \partial_1 \Gamma_2 - \partial_2 \Gamma_1
}
{ =} { \partial_1 \left(  7x^3 +xy^2-4y^5  \right) -  \partial_2 \left(   x^5-x^2y^2+3y^4   \right)
}
{ =} {  21x^2 +y^2 - 5x^4 -2xy^2
}
{ } {
}
}
{}
{}{.}

 }

__NOEDITSECTION__

\inputaufgabeklausurloesung
{0}

{

}
{  }

__NOEDITSECTION__

\inputaufgabeklausurloesung
{6}

{

Beweise den Satz über die Retraktionen auf den Rand auf Mannigfaltigkeiten.

}
{

Der Rand
\mathl{\partial M}{} ist nach
[[Mannigfaltigkeit mit Rand/Orientierung/Randorientierung/Fakt|Satz 21.8 (Differentialgeometrie (Osnabrück 2023))]]
eine
\definitionsverweis {orientierte}{}{}
\definitionsverweis {differenzierbare Mannigfaltigkeit}{}{}
\zusatzklammer {ohne Rand} {} {.}
Daher gibt es nach
[[Nullstellenfreie Volumenform/Orientierte (Karten) Mannigfaltigkeit/Äquivalenz/Fakt|Satz 22.11 (Differentialgeometrie (Osnabrück 2023))]]
eine
\definitionsverweis {stetig differenzierbare}{}{}
\definitionsverweis {positive Volumenform}{}{}
$\tau$ auf
\mathl{\partial M}{.} Es ist
\mathl{\int_{  \partial M  }  \tau >0}{.} Die
\definitionsverweis {äußere Ableitung}{}{}
der Volumenform $\tau$ ist $0$. Nehmen wir an, dass es eine
\definitionsverweis {stetig differenzierbare Abbildung}{}{}
\maabbdisp {\varphi} { M } { \partial M
} {}
mit
\mathl{\varphi \vert_{\partial M} = \operatorname{Id}_{\partial M}}{} gebe. Dann ist die
\definitionsverweis {zurückgezogene Form}{}{}

\mathl{\varphi^* \tau}{} eine
$(n-1)$-\definitionsverweis {Differentialform}{}{}
auf $M$, deren Einschränkung auf den Rand mit $\tau$ übereinstimmt. Daher gilt unter Verwendung von
[[Mannigfaltigkeit mit Rand/Satz von Stokes/Fakt|Satz 23.2 (Differentialgeometrie (Osnabrück 2023))]]
und
[[Differentialform/Mannigfaltigkeit/Äußere Ableitung/Eigenschaften/Fakt|Satz 20.4 (Differentialgeometrie (Osnabrück 2023))  (5)]]

\mavergleichskettealign
{\vergleichskettealign
{
\int_{  \partial M  }  \tau
}
{ =} {
\int_{  \partial  M  }  \varphi^* \tau
}
{ =} {
\int_{   M  }  d(\varphi^* \tau)
}
{ =} {
\int_{   M  }  \varphi^* (d\tau)
}
{ =} {
\int_{   M  }  \varphi^* (0)
}
}
{
\vergleichskettefortsetzungalign
{ =} { 0
}
{ } {}
{ } {}
{ } {}
}
{}{.}
Dies ist ein Widerspruch.

 }

 [[Kategorie:Latexseite]]
